\hypertarget{6-conclusion}{%
\section{Conclusion}\label{6-conclusion}}

This paper presented a reproducible DES+RL benchmark for operational
resilience control in a military food supply chain, grounded in
Garrido-Rios' 13-operation simulation model~\cite{garrido2017}. The
experimental design ultimately answers a benchmark-design question: does an
apparent reinforcement-learning advantage survive when the static comparator
is optimized over the same decision contract?

Under the thesis-grounded Track A contract, where the agent controls
upstream inventory and assembly shifts, dense common-random-number
static frontiers and multiple RL configurations across reward modes,
observation versions, and algorithms leave no publishable dynamic-control
advantage. This boundary is a property of the thesis-grounded decision
family, not of upstream control per se: the identified factorial shows
that a richer upstream/shift parameterization under closed-loop control
does convert headroom, so what Track A lacked was an action
parameterization whose adaptive movement matters, with the downstream
dispatch bottleneck additionally outside its controllable interface.

Under the DES-preserving Track B operational extension, which adds downstream transport
dispatch at Op10 and Op12 (an 8-dimensional action contract) while
preserving the DES backbone, PPO improves Garrido/Excel resilience over
the best of a restricted dense downstream-dispatch grid of static comparators under paired
common-random-number evaluation (+0.000438 Excel ReT across 10 training
seeds, 95\% CI entirely positive),
with matching gains in service continuity and recovery-tail metrics. A
local 3$\times$3 bound over the nearest upstream static quantity controls
at the best downstream cell narrows but does not overturn the PPO
advantage. A
5-seed identified factorial decomposes the gain: most of it comes from
closed-loop control of the upstream and shift levers the static family
holds fixed, downstream dispatch access adds a causally identified
increment, and dispatch-only control does not beat the dense frontier
that already optimizes that subspace --- so action-space size alone does
not explain the result; the
advantage holds for current and increased Garrido-native risk levels at
both evaluated horizons with fully positive seed-clustered confidence
intervals, while the severe service-floor regime is a disclosed boundary
where the frozen policy loses narrowly; and two
independent tests --- an observation-masked retrain and a
privileged-access static comparator --- rule out that the effect depends
on the learned policy exploiting simulator-only regime information
unavailable to any static or heuristic baseline. However, these controls do
not address comparator scope. In the terminal same-contract challenge, a
calibration-only constant full-contract policy attains Excel ReT 0.005906366
versus 0.005888317 for canonical PPO; PPO minus static is $-0.000018049$
with two-way CI95 $[-0.000028615,-0.000008087]$. This reverses the earlier
headline and retires the adaptive-superiority claim.

The central lesson is methodological: comparator decision rights are part of
the treatment definition. A large gain can be stable across seeds, tapes,
tail metrics, and information controls while remaining an artifact of a
static family that holds learner-controlled dimensions fixed. For researchers
and practitioners, the implication is direct: optimize a credible same-
contract non-learning policy before attributing resilience gains to
adaptation, architecture, or bottleneck access.
